\chapter{Monopoly}
\label{chap:monopoly}

\section{Learning objectives}

After completing this chapter, you should be able to:
\begin{itemize}
  \item identify the sources of monopoly power;
  \item derive a monopolist's output and price choice;
  \item compare monopoly with perfect competition; and
  \item evaluate the welfare and policy implications of monopoly.
\end{itemize}

\section{What is monopoly?}

A \keyterm{monopoly} is a market with one seller and substantial barriers to entry. A monopolist has market power: it recognises that selling more output requires a lower price. Monopoly is not defined merely by a firm's large size; it depends on whether alternatives and potential entry constrain the firm.

Sources of market power include legal protection such as patents, control of an essential input, network effects, economies of scale, and regulatory barriers. Some sources may reward innovation or reduce costs; others may arise from exclusionary practices or political privilege. Their policy implications differ.

\section{Demand, revenue, and marginal revenue}

The monopolist faces the market demand curve. If inverse demand is $p(Q)$, total revenue is
\[
TR(Q)=p(Q)Q.
\]
Unlike a competitive firm, the monopolist's marginal revenue is below price. To sell one additional unit, it must lower the price not only on that unit but normally on earlier units as well.

For linear inverse demand $p=a-bQ$,
\[
TR(Q)=aQ-bQ^2
\qquad\text{and}\qquad
MR(Q)=a-2bQ.
\]
The marginal-revenue curve has the same vertical intercept as demand and twice its slope.

\section{Profit maximisation}

The monopolist chooses output where marginal revenue equals marginal cost:
\[
MR(Q)=MC(Q).
\]
It then uses the demand curve to find the price consumers are willing to pay for that quantity. The sequence matters: a monopolist does not set price equal to marginal cost, except in special cases.

Suppose $p=50-2Q$ and $C(Q)=10+2Q$. Then $MR=50-4Q$ and $MC=2$. Setting $MR=MC$ gives $Q_M=12$. Demand then gives $p_M=26$. Profit is revenue minus cost:
\[
\pi=26(12)-(10+2\cdot12)=278.
\]

\section{Monopoly versus competition}

With the same demand and marginal cost, a competitive market produces where $p=MC$. In the example, competition yields $Q_C=24$ and $p_C=2$. Monopoly restricts output and charges a higher price. Some consumers who value the good above marginal cost do not buy it, creating a \keyterm{deadweight loss}.

Monopoly also redistributes surplus. Consumers who continue to buy pay more, transferring part of their surplus to the firm. This transfer is not itself a deadweight loss; the loss arises from units that are not traded even though their value to consumers exceeds their production cost.

\section{Price discrimination}

\keyterm{Price discrimination} occurs when a seller charges different prices to buyers for reasons other than differences in cost. It requires market power, the ability to distinguish buyers or their willingness to pay, and limited resale between buyers.

Under perfect first-degree price discrimination, the seller charges each buyer their maximum willingness to pay. Output can reach the efficient level, but the seller captures almost all surplus. More common forms include quantity discounts, student prices, peak pricing, and different fares for business and leisure travellers.

\section{Policy toward monopoly}

Possible responses include competition law, price regulation, public ownership, procurement, and patent policy. No response is universally best. A natural monopoly, where one firm can serve the market at lower cost than several firms, may require regulation rather than forced fragmentation. Innovation incentives, administrative capacity, and distributional goals must be considered alongside static efficiency.

\section*{Chapter Review}

\subsection*{Key concepts}
\begin{itemize}
  \item \textbf{Monopoly}: one seller protected by barriers to entry.
  \item \textbf{Market power}: ability to influence price or other market terms.
  \item \textbf{Marginal revenue}: additional revenue from an additional unit sold.
  \item \textbf{Deadweight loss}: surplus lost because monopoly output is below the efficient level.
  \item \textbf{Price discrimination}: charging different buyers different prices not explained by cost differences.
\end{itemize}

\section*{Exercises}

\subsection*{Conceptual questions}
\begin{enumerate}[label=\arabic*.]
  \item List four possible barriers to entry. Which may be socially beneficial in some circumstances?
  \item Why is marginal revenue below price for a monopoly with downward-sloping demand?
  \item Distinguish a transfer of surplus from deadweight loss.
  \item What conditions are needed for price discrimination?
\end{enumerate}

\subsection*{Graphical exercises}
\begin{enumerate}[resume,label=\arabic*.]
  \item Draw demand, marginal revenue, and constant marginal cost. Mark monopoly output and price.
  \item On the same diagram, mark the competitive output and the monopoly deadweight loss.
\end{enumerate}

\subsection*{Analytical exercises}
\begin{enumerate}[resume,label=\arabic*.]
  \item A monopolist faces $p=100-2Q$ and $C(Q)=10+20Q$. Find monopoly quantity, price, profit, and the competitive quantity.
  \item For Exercise 7, calculate the deadweight loss relative to competition.
  \item Explain whether a firm that charges students a lower price than professionals is necessarily behaving unfairly or inefficiently.
\end{enumerate}

\subsection*{Exam-style multiple-choice questions}
\begin{enumerate}[resume,label=\arabic*.]
  \item A profit-maximising monopolist chooses output where: (a) $p=MC$; (b) $MR=MC$; (c) $TR=TC$; or (d) $MR=p$.
  \item Relative to perfect competition, monopoly normally produces: (a) more output at a lower price; (b) less output at a higher price; (c) the same output; or (d) zero output.
  \item First-degree price discrimination can: (a) eliminate all trade; (b) increase output toward the efficient level; (c) make marginal revenue exceed price; or (d) remove market power.
\end{enumerate}

\subsection*{PPE discussion question}

\noindent When a firm has a dominant position because of a successful innovation, should government limit its market power? Discuss the trade-off between static consumer welfare and incentives for innovation.
